\section{Component Replays and Additional Evidence}\label{app:integration}
\begin{table}[htbp]\centering\small
\begin{tabular}{llrrrrr}
\toprule Seed & Arm & Learned & Exact CMT & Exact P & Exact CAND & Exact KW \\
\midrule
\closureinput tables/integration-components.tex
\bottomrule
\end{tabular}
\caption{All six selected BE models on their original 1,000-case batteries.
CMT replaces only confirmation, P only participation, CAND only candidate;
the dependent learned confirmation is recomputed when messages change.
Exact KW replaces all three. Linial remains learned. These are 30,000
diagnostic executions on consumed cases, not 30 training replicas.}
\label{tab:components-full}
\end{table}

All baselines reproduce complete stored terminal fields in their original
cohort 32/inference-batch 1,024 float32 CPU context. The six checkpoints are
chosen because their three seeds include every functional failure in the
input-context study; the two nonfailing native arms are paired controls.
No checkpoint is retrained, no failing case removed and no new primary
endpoint introduced. Baseline, COMMIT-only, participation-only,
candidate-only and all-KW modes are fixed before outcomes. All six models
and all cases remain in the output, including ineffective interventions.
The selected CF failure count is 2,595, which is also the complete CF
failure inventory for that study's original 10,000 cases. Rescue is a
statement about those cases, not a population success probability.

\paragraph{Contract rejection need not mean an invalid write.}
The baseline diagnostic records 64 and 20 active acceptances rejected by the
contract in native seeds 105001 and 105010. A subsequent unchanged-policy
observer reproduces both 1,000-case batteries exactly and checks the actual
KW decoding. Every one of these 84 writes reproduces its old auxiliary
color and satisfies the new bound. The contexts $(B,a,c,D)$ are
$(8,6,2,6)$, $(12,9,1,9)$ and $(14,11,3,11)$, with respective counts
53/15, 6/1 and 5/4 in the two seeds. They are errors against the specified
contract but do not alter that transition, as Proposition~\ref{prop:relevance}
predicts. The frozen counter named \code{learned_unsafe_accepts} in this BE
diagnostic means active contract-rejected acceptances, not observed
collisions; it must not be conflated with transfer's forbidden-write counter.
This observer is posthoc and does not revise the frozen intervention outputs.

\paragraph{Campaigns and label budgets.}
The historical diagnostic uses five seeds, unequal acquisition budgets and
development-reused queries. The 40-model input-context study uses ten paired
seeds, a shared 3,000-target base, 3,600 auxiliary targets per arm,
1,500 fixed updates and a fine-key hash holdout. The 100-model label-budget
study uses five paired seeds, 4,000 updates and different intervention
families, support and endpoints; its 1,500-update snapshots are secondary.
The five-model neighbor-message follow-up freezes its modified acquisition
after inspecting that study and uses a fresh v2 evaluation battery. The
studies are not pooled and their equal-type metrics differ: the former
averages accuracy equally over types, while the latter also class-balances
its type endpoint. Numerical bounds used by a generator are not necessarily
observed by the COMMIT head.

The label-budget study has 32/40 fully passing BE checkpoints; only two of
five high-CF seeds meet the 99\% unseen-type criterion. All acquisition
variants are retained, including the no-force and late-bound conditions.
Complete component-group accounting gives slightly different native/CF label
totals; neither identical logical support nor equivalent class information
follows from a matched nominal budget.

\paragraph{What further experiments would identify.}
Attributing an advantage specifically to counterfactual label generation
would require a control with informative positive and negative
labels, matched support/balancing, updates and checkpoint selection.
Generalization to unseen interventions would require holding out complete
families and checking novelty at the effective model input, rather than only
at a generator parameter. Neither stronger claim is made here.
